\section{Conclusion}

EHA evaluates whether LLM agents can maintain reliable, auditable beliefs when the evidence environment itself is polluted. The opaque frontier cohort run shows why this matters. Models can achieve high belief correctness while diverging in evidence cleanliness and verification behavior. Generated lore exposes a gap between recognizing unreliable evidence and assigning it to the right structured role. Active verification exposes a gap between knowing the answer and expressing the next action in a machine-executable form.

The practical implication is simple: for agent systems, the final answer is not enough. Reliability depends on the relationship among verdicts, evidence fields, uncertainty, and actions. EHA provides a controlled diagnostic way to measure that relationship in a setting where provenance and pollution are auditable.

The next step is not to treat this release as a finished universal benchmark, but to scale the task set, stress-test surface cues, broaden schema ablations, and complete independent human audit of active-verification judgments. Those extensions would test whether the same decomposition remains stable as EHA moves from a mechanism study toward a larger benchmark.
